%% EE568 -- Design of Electrical Machines
%% Course lecture notes, typeset from handwritten notes (Spring 2026)
%% Compile: pdflatex EE568_Lecture_Notes.tex  (run twice for the ToC)

\documentclass[12pt,a4paper]{article}

\usepackage[margin=2.5cm]{geometry}
\usepackage{amsmath, amssymb}
\usepackage{graphicx}
\usepackage{float}
\usepackage{hyperref}
\usepackage{enumitem}
\usepackage{booktabs}
\usepackage{array}
\usepackage{fancyhdr}
\usepackage[table]{xcolor}
\usepackage[T1]{fontenc}
\usepackage[utf8]{inputenc}

\setlength{\headheight}{15pt}
\setlength{\parskip}{6pt}

\newcommand{\unit}[1]{\ensuremath{\,\mathrm{#1}}}
\newcommand{\degree}{\ensuremath{^{\circ}}}

\pagestyle{fancy}
\fancyhf{}
\lhead{EE568 -- Lecture Notes}
\rhead{Design of Electrical Machines}
\rfoot{\thepage}

\hypersetup{colorlinks=true, linkcolor=blue!70!black, urlcolor=blue!70!black}

\title{\textbf{EE568 -- Design of Electrical Machines}\\[4pt]
       \large Lecture Notes}
\author{Middle East Technical University \\
        Department of Electrical and Electronics Engineering \\[4pt]
        Compiled by Sefer \.{I}DACI}
\date{Spring 2026}

\begin{document}
\maketitle
\thispagestyle{empty}

\begin{abstract}
\noindent
These notes were typeset from the handwritten course notes of EE568
``Design of Electrical Machines'' at METU, Spring 2026. They follow the
flow of the lectures: starting from magnetic-circuit fundamentals, then
building up to the design of distributed and fractional-slot windings,
to the output equation that sizes a machine, and finally to the
practical constraints (mechanical, magnetic, thermal) that shape the
final geometry. The classical references are the books by
Pyrh\"onen et~al.\ (\textit{Design of Rotating Electrical Machines},
Wiley) and Hanselman (\textit{Brushless Permanent Magnet Motor Design},
Magna Physics); equation numbering does not follow either book.
\end{abstract}

\tableofcontents
\newpage

% =====================================================================
\section{Magnetic Fundamentals}
\label{sec:magnetics}

\subsection{Electric--Magnetic Analogy}

A magnetic circuit obeys the same network laws as an electric circuit
once the analogous quantities are identified
(Table~\ref{tab:analogy}).
The driving quantity for current in the electric circuit is the voltage
$V$; in the magnetic circuit it is the magnetomotive force (MMF)
$NI$, the product of the coil turns and the current they carry.
The opposing quantity is the resistance $R$ on the electric side and
the reluctance $\mathcal{R}$ on the magnetic side.
The resulting flow quantities are current $I$ and flux $\Phi$:
\begin{equation}
  V = IR
  \qquad\Longleftrightarrow\qquad
  NI = \Phi\,\mathcal{R}
\end{equation}

\begin{table}[H]
  \centering
  \caption{Analogous quantities in electric and magnetic circuits.}
  \label{tab:analogy}
  \begin{tabular}{lll}
    \toprule
    Quantity (electric) & (magnetic) & Definition \\
    \midrule
    Voltage $V$\unit{[V]}             & MMF $NI$\unit{[A\!\cdot\!t]}   & driving source \\
    Current $I$\unit{[A]}             & Flux $\Phi$\unit{[Wb]}         & flow quantity \\
    Current density $J = \sigma E$\unit{[A/m^2]} & Flux density $B = \mu H$\unit{[Wb/m^2 = T]} & flow per area \\
    Field $E$\unit{[V/m]}             & Field $H$\unit{[A/m]}          & gradient of source \\
    Conductivity $\sigma$             & Permeability $\mu$              & material property \\
    Resistance $R = \ell/(\sigma A)$  & Reluctance $\mathcal{R} = \ell/(\mu A)$ & opposes flow \\
    \bottomrule
  \end{tabular}
\end{table}

For a series magnetic path of $n$ segments,
$\mathcal{R}_\text{tot} = \sum_i \ell_i/(\mu_i A_i)$.
The inductance of an $N$-turn coil wound on a magnetic circuit of total
reluctance $\mathcal{R}$ is
\begin{equation}
  L = \frac{\lambda}{I} = \frac{N\Phi}{I} = \frac{N \cdot NI/\mathcal{R}}{I}
    = \frac{N^2}{\mathcal{R}}
  \label{eq:L_rel}
\end{equation}

\subsection{Air is a Circuit Element}

A magnetic circuit cannot be drawn the way an electric circuit can,
because flux always closes on itself \emph{including through air}.
Whenever the path through a high-permeability core encounters a gap,
the air gap appears in series with a large reluctance
$\mathcal{R}_g = g/(\mu_0 A)$.
Since $\mu_\text{Fe}/\mu_0$ is typically $10^3$--$10^4$, even a
millimetre-thick gap dominates the total reluctance of a machine
magnetic circuit. This is the central reason why electrical machines
are sized by their air-gap geometry rather than by their iron
geometry.

\subsection{B--H Curves and Permeability}

Magnetic materials are non-linear: $B(H)$ saturates and follows a
hysteresis loop on reversal.
Two different ``permeabilities'' are commonly quoted
(Figure~\ref{fig:bh}):
\begin{itemize}[noitemsep]
  \item \textbf{Apparent permeability} $\mu = B/H$ -- the secant slope
        from the origin to the operating point. Used when the material
        is treated as linear at one operating point.
  \item \textbf{Incremental permeability} $\mu = dB/dH$ -- the local
        tangent slope. Relevant for small-signal AC excitation
        superimposed on a DC bias.
\end{itemize}

\begin{figure}[H]
  \centering
  \begin{minipage}{0.45\textwidth}\centering
    \textit{B-H curve, single-valued} \\[2pt]
    (linear-material approximation; secant and tangent slopes shown).
  \end{minipage}\hfill
  \begin{minipage}{0.45\textwidth}\centering
    \textit{Hysteresis loop} \\[2pt]
    (coercive force $H_c$; loop area = hysteresis loss per cycle).
  \end{minipage}
  \caption{Magnetic material behaviour. Hysteresis loss is independent
           of frequency per cycle, so the total loss scales linearly
           with frequency.}
  \label{fig:bh}
\end{figure}

\paragraph{Coercive force and hard vs soft materials.}
The coercive force $H_c$ is the value of $H$ required to drive $B$ back
to zero from remanence. Materials with small $H_c$ are easily
demagnetised and are called \emph{soft} (used for cores); materials
with large $H_c$ are hard to demagnetise and are called \emph{hard}
(used for permanent magnets). The product $|B||H|$ in the second
quadrant defines the \emph{maximum energy product} -- the figure of
merit for a permanent magnet (Section~\ref{sec:pm}).

\subsection{Core Losses}

Two mechanisms convert magnetic flux variation into heat in
ferromagnetic cores:
\begin{enumerate}[noitemsep]
  \item \textbf{Hysteresis loss.} Loop area per cycle; loss per second
        $\propto f$ (independent of the magnetic properties per cycle,
        related to the cyclic reversal of domain orientation).
        For a 50\unit{Hz} machine the loop traverses 50 times per
        second.
  \item \textbf{Eddy-current loss.} Induced EMFs in the laminations
        circulate currents that dissipate in the core resistance.
\end{enumerate}

\paragraph{Eddy-current loss derivation.}
A small loop of resistance $R$ in the laminating direction sees the
induced EMF $e = N\,d\Phi/dt$.
If the flux is sinusoidal with peak $\hat{B}$ and frequency $f$:
$e \propto \hat{B}\,f$, so $V \propto \hat{B}\,f$ and
$P_\text{eddy} = V^2/R \propto \hat{B}^2 f^2$.
Including the lamination conductivity $\sigma$ and thickness
explicitly,
\begin{equation}
  P_\text{eddy} \;\propto\; \sigma\,\hat{B}^2\,f^2
  \label{eq:peddy}
\end{equation}
Lamination of the core breaks the loop into many small parallel
sheets, raising the effective $R$ and reducing $P_\text{eddy}$.

\paragraph{Total core loss.}
$P_\text{core} = P_\text{hyst} + P_\text{eddy}$.
At 50\unit{Hz} this is dominated by hysteresis; at high frequency
(PWM-driven motors) eddy losses dominate.

\subsection{Faraday's Law and Inductance}

The induced voltage on an $N$-turn coil is
\begin{equation}
  e = -\,N\frac{d\Phi}{dt}
    = -\,\frac{d\lambda}{dt}
    = -\,\frac{d(N\Phi)}{dt}
\end{equation}
where $\lambda = N\Phi$ is the flux linkage.
The inductance is defined as
\begin{equation}
  L = \frac{d\lambda}{dI}
    = N\frac{d\Phi}{dI}
\end{equation}
For a \emph{linear} material with constant $\mathcal{R}$,
$\Phi = NI/\mathcal{R}$ so $L = N^2/\mathcal{R}$ as in
Eq.~\eqref{eq:L_rel}. For a non-linear material with saturation,
$\lambda(I)$ is a curve that bends over (Figure~\ref{fig:lambda_I});
$L = \lambda/I$ is then \emph{not} a constant and the value
$\tfrac{1}{2}LI^2$ is \emph{not} the stored energy.

\begin{figure}[H]
  \centering
  \textit{$\lambda$--$I$ curve} (linear region near origin, saturating at high $I$;
   $L=\lambda/I$ only valid in the linear region).
  \caption{Flux linkage as a function of current.}
  \label{fig:lambda_I}
\end{figure}

\subsection{Non-Uniform Flux, Fringing and Leakage}

Real magnetic circuits depart from the ideal in several ways:
\begin{itemize}[noitemsep]
  \item \textbf{Fringing flux} bulges outward in the air gap; the
        effective gap cross-section is larger than the iron face.
  \item \textbf{Leakage flux} bypasses the intended path (slot-leakage,
        end-winding leakage, magnet-to-magnet leakage).
  \item Non-uniform flux density across a section can arise because the
        inner part of a curved iron path is shorter and thus has lower
        reluctance, attracting more flux ($B_\text{in} > B_\text{out}$).
\end{itemize}
These effects are accounted for either by Carter's coefficient
(Section~\ref{sec:carter}) or by an explicit reluctance network.

\subsection{Energy Storage in the Magnetic Field}

For a linear material the energy density is
\begin{equation}
  w = \int H\,dB = \frac{1}{2}\,B H = \frac{B^2}{2\mu}
  \label{eq:wdensity}
\end{equation}
and the total stored energy in a volume $V$ is
\begin{equation}
  W = \tfrac{1}{2}\,L I^2
    = \frac{1}{2}\,\frac{B^2}{\mu}\,V
\end{equation}
At fixed $B$, low $\mu$ stores more energy per unit volume; this is
why an air gap is added when a large amount of energy must be stored
(inductors, transformers in flyback mode, etc.). Conversely, for
\emph{zero-energy} flux paths (the iron in a machine) high
$\mu$ minimises the stored field energy and the magnetising current.

\paragraph{Maxwell stress.}
The tangential stress on an iron surface in a normal field $B$ is
\begin{equation}
  \sigma_\text{Maxwell} = \frac{B^2}{2\mu_0}
  \quad\unit{[N/m^2 = Pa]}
  \label{eq:maxwell}
\end{equation}
This is the same as the energy density of Eq.~\eqref{eq:wdensity};
the units check because $\unit{J/m^3} = \unit{N/m^2}$. The integral of
this stress over the air-gap surface times the rotor radius gives the
electromagnetic torque (Section~\ref{sec:sizing}).

% =====================================================================
\section{AC Machine Equivalent Circuit}
\label{sec:eqckt}

The per-phase steady-state equivalent circuit of an AC machine
(induction motor with the slip-dependent rotor branch as an example)
is shown in Figure~\ref{fig:eqckt}.
The stator branch has $R_s + jX_s$, where $X_s = \omega L_{s\sigma}$
is the stator leakage reactance.
The shunt branch carries the magnetising current and the iron losses:
$R_c$ in parallel with $jX_m = j\omega L_m$.
The rotor branch (referred to the stator) is $R_r' + jX_r'$, with the
``mechanical resistance'' $R_r'(1-s)/s$ representing the converted
mechanical power.

\begin{figure}[H]
  \centering
  \textit{Per-phase equivalent circuit: $X_s, R_s$ -- $\;[R_c\parallel jX_m]\;$ -- $X_r', R_r', \, R_r'(1-s)/s$.}
  \caption{Per-phase steady-state equivalent circuit of an induction
           machine.}
  \label{fig:eqckt}
\end{figure}

\paragraph{Leakage inductance vs.\ magnetising inductance.}
The leakage path is partly air, so $L_{s\sigma}$ is closer to a linear
inductor (small saturation) and is dominated by the slot/end-winding
geometry rather than the iron. The magnetising inductance $L_m$ is set
by the air gap and the iron and saturates with $I_m$.

% =====================================================================
\section{Integral-Slot Distributed Windings}
\label{sec:integral}

A polyphase winding placed in slots produces an MMF wave that travels
around the air gap. The quality of this wave -- how close it is to a
pure sinusoid, how large the fundamental component is relative to the
peak ampere-conductor count -- is described by the
\emph{winding factor} $k_w$.

\subsection{Slots per Pole per Phase}

For a machine with $N_s$ stator slots, $m$ phases and $2p$ poles,
\begin{equation}
  q = \frac{N_s}{2p\,m}
  \label{eq:q}
\end{equation}
is the number of slots per pole per phase. When $q$ is a positive
integer, the winding is \emph{integral-slot} (this section); when $q$
is a non-integer rational, the winding is \emph{fractional-slot}
(Section~\ref{sec:fractional}).

The slot pitch in electrical degrees is
\begin{equation}
  \alpha_e = \frac{p \cdot 360\degree}{N_s}
\end{equation}

\paragraph{Worked example -- 6-slot, 2-pole, 3-phase machine.}
Here $p=1$, $m=3$, $N_s=6$, so $q=1$ and $\alpha_e = 60\degree$.
The phase assignment to the six slots is
\begin{center}
  \begin{tabular}{c|cccccc}
    \toprule
    Slot   & 1 & 2 & 3 & 4 & 5 & 6 \\
    Phase  & A$+$ & C$-$ & B$+$ & A$-$ & C$+$ & B$-$ \\
    \bottomrule
  \end{tabular}
\end{center}
which yields a 60\degree-sector pattern repeating every pole-pair.

\subsection{MMF of a Concentrated Coil}

A single-slot full-pitch coil produces a rectangular MMF wave of
amplitude $NI/2$ on each side of the magnetic axis, switching sign at
$\pm\pi$ (electrical).
The Fourier expansion of this square wave is
\begin{equation}
  \mathrm{MMF}(\theta) = \frac{4}{\pi}\cdot\frac{NI}{2}
  \sum_{n=1,3,5,\ldots}\frac{1}{n}\sin(n\theta)
\end{equation}
so the fundamental amplitude is
\begin{equation}
  \widehat{\mathrm{MMF}}_1 = \frac{4}{\pi}\cdot\frac{NI}{2}
\end{equation}
A \emph{distributed} winding (multiple coils per phase per pole spread
across $q$ adjacent slots) produces a stepped waveform that approaches
a sinusoid as $q$ grows, attenuating the high-order harmonics.

\subsection{Distribution Factor $k_d$}

When $q$ coils share one pole-phase belt, their individual EMF
phasors are offset by $\alpha_e$ each, so the resultant is the vector
sum (a chord of the circumscribed polygon) rather than the arithmetic
sum:
\begin{equation}
  k_d = \frac{\text{vector sum of }q\text{ coil EMFs}}
             {\text{arithmetic sum}}
\end{equation}
The geometric-series identity gives
\begin{equation}
  k_{d,n} = \frac{\sin\!\bigl(n\,q\,\alpha_e/2\bigr)}
                 {q\,\sin\!\bigl(n\,\alpha_e/2\bigr)}
  \label{eq:kd}
\end{equation}
For the fundamental ($n=1$), $k_d$ is just below unity and increases
toward 1 as $q$ increases (smoother distribution). For higher
harmonics, $k_d$ drops sharply, which is the harmonic-attenuation
benefit of distribution.

\subsection{Pitch Factor $k_p$}

A \emph{full-pitch} coil has its two sides exactly one pole pitch
apart, so the two side EMFs are 180\degree\ out of phase and add
constructively in the coil loop.
A \emph{short-pitch} coil has its sides closer than one pole pitch by
some fraction $\lambda$ of the pole pitch, $0 < \lambda \le \pi$
(electrical). The vector subtraction of the two side EMFs has
magnitude $2\hat{E}\sin(\lambda/2)$ instead of $2\hat{E}$:
\begin{equation}
  k_{p,n} = \sin\!\left(\frac{n\,\lambda}{2}\right)
  \label{eq:kp}
\end{equation}
For the fundamental, $k_p \le 1$.
At $\lambda = \pi(1 - 1/n)$, $k_{p,n} = 0$ and the $n$-th harmonic is
eliminated entirely.
Common choices:
\begin{itemize}[noitemsep]
  \item $\lambda = 5/6\,\pi$ (150\degree, ``5/6 pitched'') --
        suppresses the 11th and 13th in a three-phase machine.
  \item $\lambda = 2/3\,\pi$ (120\degree) -- eliminates the third
        harmonic (and triplens).
  \item $\lambda = 4/5\,\pi$ (144\degree) -- eliminates the fifth.
\end{itemize}

\paragraph{Why eliminate harmonics.}
Space harmonics of the MMF do not contribute average torque (only the
fundamental does), but they produce iron losses, torque ripple,
acoustic noise, and parasitic asynchronous torques. Even harmonics
imply DC bias in the field distribution and are absent by construction
in a balanced winding (each half of the loop sees both $+$ and $-$
poles).

\subsection{Winding Factor $k_w$}

The total winding factor combines both effects:
\begin{equation}
  k_{w,n} = k_{d,n}\cdot k_{p,n}
  \label{eq:kw}
\end{equation}
A high $k_{w,1}$ maximises torque density; a low $k_{w,n}$
for $n>1$ reduces harmonic content.

\subsection{Worked Example -- 180-slot, 50 Hz, 10-Pole Alternator}

A rotating field produced by a balanced 3-phase 50\unit{Hz}, 600\unit{rpm}
alternator has the magnetic flux distribution
\begin{equation}
  B(\theta) = 0.9\sin\theta + 0.25\sin 3\theta + 0.18\sin 5\theta
\end{equation}
The machine has 180 slots, double-layer Y-connected winding, each coil
of 3 turns and coil span 15 slots. Armature diameter $D = 1.35\unit{m}$,
effective axial length $L = 0.5\unit{m}$.

\emph{Number of poles.}
$f_m = 600/60 = 10\unit{Hz}$, $f_e = 50\unit{Hz}$, so
$p = f_e/f_m = 5$ pole-pairs (10 poles).

\emph{Slots per pole per phase.}
$q = 180/(3\cdot 10) = 6$.

\emph{Pole area.}
\begin{equation}
  A_\text{pole} = \frac{\pi D L}{2p}
                = \frac{\pi \cdot 1.35 \cdot 0.5}{10}
                = 0.212\unit{m^2}
\end{equation}

\emph{Slot pitch (electrical) and coil pitch.}
$\alpha_e = (360\degree/180)\cdot p = 10\degree$ electrical;
coil span 15 slots $\Rightarrow$ $\lambda = 150\degree$ electrical.

\emph{Distribution factors.}
\begin{align*}
  k_{d,1} &= \sin(6\cdot 10\degree/2)/(6\sin(10\degree/2))
           = \sin 30\degree / (6\sin 5\degree)
           = 0.956\\
  k_{d,3} &= \sin 90\degree / (6\sin 15\degree) = 0.643\\
  k_{d,5} &= \sin 150\degree / (6\sin 25\degree) = 0.197
\end{align*}

\emph{Pitch factors.}
\begin{align*}
  k_{p,1} &= \sin(150\degree/2) = \sin 75\degree = 0.966\\
  k_{p,3} &= \sin 225\degree = -0.707\\
  k_{p,5} &= \sin 375\degree = \sin 15\degree = 0.259
\end{align*}

\emph{Winding factors.}
\begin{equation}
  k_{w,1} = 0.923,\quad
  k_{w,3} = -0.455,\quad
  k_{w,5} = 0.051
\end{equation}

\emph{Peak flux per pole (fundamental).}
\begin{equation}
  \hat{\Phi}_1 = \frac{2}{\pi}\,\hat{B}_1\,A_\text{pole}
               = \frac{2}{\pi}\cdot 0.9\cdot 0.212 = 0.1214\unit{Wb}
               = 121.4\unit{mWb}
\end{equation}

\emph{Turns per phase.}
For $q=6$, $2p=10$, 3 turns per coil, double-layer:
$N_s = q\cdot 2p\cdot 3 = 180$ turns per phase.

\emph{Induced phase voltage (per harmonic) by Faraday.}
The classical alternator formula is
\begin{equation}
  V_\text{ind,phase}
    = \frac{2\pi}{\sqrt{2}}\,f\,N_s\,\hat{B}\,A_\text{pole}\,k_{w,n}
    = 4.44\,f\,N_s\,\hat{B}\,A_\text{pole}\,k_{w,n}
\end{equation}
For the three harmonics:
\begin{center}
  $V_1 \approx 4480\unit{V}$,\quad $V_3 \approx -613\unit{V}$,\quad $V_5 \approx 49.5\unit{V}$ (rms)
\end{center}
The phase voltage is the RMS sum
$V_\text{ph} = \sqrt{V_1^2 + V_3^2 + V_5^2}$.
The Y-connection eliminates the third harmonic from the line-to-line
voltage (triplens are in phase across the three phases), so the
\emph{line} voltage waveform is much closer to sinusoidal than the
phase voltage.

% =====================================================================
\section{Fractional-Slot Windings}
\label{sec:fractional}

A \emph{fractional-slot} winding has $q \notin \mathbb{Z}$.
Such windings are common in low-speed, high-pole-count machines
(direct-drive PM motors, in particular). Their main practical
advantages are:
\begin{itemize}[noitemsep]
  \item Short end-windings (when $q<1$, tooth-coil), lower copper
        volume and lower DC resistance.
  \item Reduction of cogging torque, because the slot/pole
        symmetry is broken.
  \item Better stator-winding fault tolerance (faulty coils can be
        isolated without unbalancing the rotating field).
\end{itemize}
Their main disadvantages are:
\begin{itemize}[noitemsep]
  \item Lower fundamental winding factor (in general) than a
        comparable distributed winding.
  \item Rich harmonic content of the MMF, including \emph{sub}-harmonics
        that can saturate the rotor back-iron of PM machines and
        cause additional rotor losses.
  \item More complex analysis: closed-form $k_d$, $k_p$ no longer
        apply; the star-of-slots method is needed instead.
\end{itemize}

\subsection{Star of Slots}

Each slot $k$ is assigned an EMF phasor at electrical angle
\begin{equation}
  \varphi_k = (k-1)\,\alpha_s \pmod{360\degree},
  \qquad
  \alpha_s = \frac{p\cdot 360\degree}{N_s}
  \label{eq:star_of_slots}
\end{equation}
The $N_s$ phasors form a star (or a regular polygon traced multiple
times, depending on $\gcd(N_s,p)$).
For a balanced three-phase winding the $360\degree$ circle is divided
into six $60\degree$ sectors:
\begin{center}
  A$+$: $[0\degree,60\degree)$;\quad
  C$-$: $[60\degree,120\degree)$;\quad
  B$+$: $[120\degree,180\degree)$;\quad
  A$-$: $[180\degree,240\degree)$;\quad
  C$+$: $[240\degree,300\degree)$;\quad
  B$-$: $[300\degree,360\degree)$
\end{center}
Each slot is assigned to the phase whose sector contains its phasor.
The $n$-th harmonic winding factor is then computed by the
phasor sum
\begin{equation}
  k_{w,n} = \frac{1}{N_{ph}}
  \left|\sum_{k\in\text{ph}^+} e^{jn\varphi_k}
        + \sum_{k\in\text{ph}^-} e^{j(n\varphi_k + \pi)}\right|
\end{equation}
where $N_{ph}$ is the number of slots per phase and ph$^+$/ph$^-$ are
the positive and negative slots of that phase.

\subsection{Balancing Condition}

For a three-phase winding to be balanced, the number of slots per
phase $N_s/m$ must be an integer:
\begin{equation}
  \frac{N_s}{m} \in \mathbb{Z}^+
  \label{eq:balance}
\end{equation}
This is necessary but not sufficient; in addition,
$\gcd(N_s, 2p)$ should be such that the phasor stars of the three
phases are identical apart from a 120\degree\ rotation.
For example:
\begin{itemize}[noitemsep]
  \item 10-pole, 30 slots: $q=1$, integral, balanced.
  \item 10-pole, 60 slots: $q=2$, integral, balanced.
  \item 10-pole, 42 slots: $q=7/5$, fractional;
        $42/3 = 14$ slots/phase is integer, balanced.
  \item 10-pole, 44 slots: $44/3$ is not integer -- not balanced.
\end{itemize}

\subsection{Slot/Pole Combinations}

Some slot/pole combinations have particularly high fundamental winding
factor and low cogging torque; tabulated examples (see Pyrh\"onen
Table~2.5) include
\begin{itemize}[noitemsep]
  \item 12-slot, 10-pole ($k_{w,1}\approx 0.933$);
  \item 12-slot, 14-pole ($k_{w,1}\approx 0.933$);
  \item 15-slot, 14-pole;
  \item 18-slot, 16-pole.
\end{itemize}
Multiples of these basic units (24/20, 24/28, \ldots) are also widely
used.

\subsection{Worked Example -- 12 Slots, 10 Poles, Single Layer}

$\alpha_s = (360\degree/12)\cdot 5 = 150\degree$ electrical.
Drawing the 12 phasors on the star and grouping them by 60\degree\
sector yields 4 slots per phase (overlapping pairs at the same
angle). The fundamental winding factor obtained from the phasor sum
is $k_{w,1}\approx 0.933$. This combination uses tooth-coil (concentrated)
windings -- each coil is wound around one stator tooth, dramatically
shortening the end winding.

% =====================================================================
\section{Machine Sizing -- The Output Equation}
\label{sec:sizing}

The classical \emph{sizing equation} relates the apparent power
delivered by an AC machine to its airgap volume, the electric loading,
the magnetic loading, and the synchronous speed.

\subsection{Magnetic and Electric Loading}

\paragraph{Specific magnetic loading $\bar B$.}
The average air-gap flux density over one pole (one pole pitch on the
bore circumference):
\begin{equation}
  \bar{B} = \frac{p\,\Phi_p}{\pi D L}
  \qquad\unit{[T]}
\end{equation}
where $\Phi_p$ is the flux per pole, $D$ the bore diameter and $L$ the
axial length. Typical values: $0.6$--$0.9\unit{T}$ for industrial
induction motors, $0.8$--$1.0\unit{T}$ for PM motors.

\paragraph{Specific electric loading $\bar A$.}
The RMS ampere-turns per unit length of air-gap circumference:
\begin{equation}
  \bar A = \frac{N_\text{turn,slot}\,I\,Q}{\pi D_i}
  \qquad\unit{[A/m]}
\end{equation}
where $Q = N_s$ is the slot count, $N_\text{turn,slot}$ the turns per
slot, $I$ the RMS phase current. Typical values: $20$--$40\unit{kA/m}$
for industrial motors, up to $80$--$100\unit{kA/m}$ for forced-cooled
machines. Higher $\bar A$ raises torque density but worsens
$I^2R$ losses and demands better cooling.

\subsection{Output Equation}

The apparent power of an $m$-phase machine is
$S = m\,E_m I_s$ where $E_m$ is the phase RMS EMF and $I_s$ the RMS
current. Substituting the alternator formula
$E_m = (2\pi/\sqrt{2})\,f\,N_s\,k_{w,1}\,\hat{B}\,A_\text{pole}/p$
and grouping with the geometric variables:
\begin{equation}
  S_i
    = \frac{\pi^2}{\sqrt{2}}\,n_\text{syn}\,k_{w,1}\,\bar A\,\hat B\,D^2 L
  \label{eq:Si}
\end{equation}
The result is usually written
\begin{equation}
  \boxed{\,S_i = C\,D^2 L\,n_\text{syn}\,}
  \qquad
  C = \frac{\pi^2}{\sqrt{2}}\,k_{w,1}\,\bar A\,\hat B
  \label{eq:C}
\end{equation}
where $C$ is the \emph{specific machine constant}, a single figure of
merit that captures how well the machine uses its volume.
$C$ depends only on the materials and the cooling capability; it lets
machines of different sizes be compared at a glance.

\paragraph{How to increase power.}
From Eq.~\eqref{eq:Si}, the three independent levers are:
\begin{enumerate}[noitemsep]
  \item Increase the volume ($D^2 L$): build a physically larger machine.
  \item Increase the loadings ($\bar A$, $\bar B$): better materials,
        better cooling, higher current density.
  \item Increase $n_\text{syn}$: spin it faster.
\end{enumerate}
The number of poles by itself does \emph{not} change $S_i$; for a
fixed $n_\text{syn}$, more poles means higher $f$ but the same output
power.

\subsection{Worked Example -- Pyrh\"onen 6.1: 4 kW Induction Motor}

Given $D=98\unit{mm}$, $L=82\unit{mm}$, $P = 4\unit{kW}$,
$f=50\unit{Hz}$, $2p=2$, rated slip $s=0.01$.

Synchronous speed $n_\text{syn} = f/p = 50\unit{Hz} = 3000\unit{rpm}$,
mechanical $\omega_\text{mech} = 2\pi f(1-s) = 99\pi\unit{rad/s}$.

Torque:
\begin{equation}
  T = \frac{P}{\omega_\text{mech}} = \frac{4000}{99\pi} \approx 12.86\unit{Nm}
\end{equation}

Air-gap (rotor surface) area:
$S_r = \pi D L = \pi \cdot 0.098 \cdot 0.082 = 0.0252\unit{m^2}$.

Tangential force and shear stress:
\begin{equation}
  F_\text{tan} = T/r = 12.86/0.049 = 262\unit{N},
  \quad
  \sigma = F_\text{tan}/S_r = 262/0.0252 \approx 10.4\unit{kPa}
\end{equation}
This is in the standard 5--35\unit{kPa} band for industrial
non-PM motors.

Specific machine constant (mechanical):
\begin{equation}
  C_\text{mech} = \frac{P}{D^2 L f_\text{sync}}
                = \frac{4000}{0.098^2\cdot 0.082\cdot 50}
                \approx 102\unit{kW\!\cdot\!s/m^3}
\end{equation}

\subsection{Aspect Ratio}

The ratio $\chi = L/D$ is a free design parameter constrained mostly
by mechanical and thermal considerations. For asynchronous machines a
common rule of thumb is
\begin{equation}
  \chi \approx \frac{\pi}{2p}\,\sqrt[3]{p}
\end{equation}
with $0.5 \lesssim D/L \lesssim 2.5$ as a typical industrial band.
Synchronous machines often have $\chi < 1$ (pancake) for direct-drive
applications and $\chi > 1$ (slender) for high-speed applications.

\subsection{Worked Example -- 30 kW Squirrel-Cage Induction Generator}

A 30\unit{kW}, 690\unit{V}, 50\unit{Hz}, 4-pole, 3-phase squirrel-cage
induction generator is to be sized.
Choose $C_\text{mech} = 150\unit{kW/m^3}$ (a typical value for a
totally-enclosed fan-cooled machine).

Aspect ratio:
\begin{equation}
  \chi \approx \frac{\pi\sqrt[3]{p}}{2p}
        = \frac{\pi\sqrt[3]{2}}{2\cdot 2}
        = 0.989 \approx 1
\end{equation}
i.e.\ $L \approx D$, so $D^2 L \approx D^3$.

Bore diameter:
\begin{equation}
  D = \sqrt[3]{\frac{P}{C_\text{mech}\,f_\text{sync}}}
    = \sqrt[3]{\frac{30\,000}{150\,000\cdot 25}}
    \approx 200\unit{mm}
\end{equation}
(taking $f_\text{sync} = f/p = 25$ rev/s for the 4-pole machine).
With $L \approx D = 200\unit{mm}$.

Mechanical speed $n_\text{mech} \approx 1474\unit{rpm}$
(slip $\sim 1.7\%$), so the torque is
$T = P/\omega = 30\,000/(2\pi\cdot 1474/60) \approx 194.4\unit{Nm}$.

Tangential stress at the rotor surface:
\begin{equation}
  \sigma_\text{tan}
    = \frac{T}{S_r\,r}
    = \frac{T}{\pi D L \cdot D/2}
    = \frac{194.4}{\pi(0.2)^3/2}
    \approx 15.5\unit{kPa}
\end{equation}
again within the standard industrial band.

Stator circumference $\pi D = 628\unit{mm}$.
For $q=2$: $N_s = 2\cdot 4\cdot 3 = 24$ slots; assuming equal slot and
tooth widths, the tooth thickness is
$\tau_\text{teeth} = 628/(2\cdot 24) = 13\unit{mm}$.
For $q=4$: $N_s=48$, $\tau_\text{teeth} = 6.5\unit{mm}$.
For $q=3$: $N_s=36$, $\tau_\text{teeth} = 8\unit{mm}$ -- a reasonable
compromise.
Coil span: about 9 slots.

% =====================================================================
\section{Mechanical Constraints}
\label{sec:mech}

\subsection{Air-Gap Selection}

The air gap $g$ is set by mechanical clearance, eccentricity tolerance,
and the level of additional losses one is willing to accept from
slot-harmonic field penetration. An empirical rule
(Pyrh\"onen Eq.~6.7) for industrial induction motors:
\begin{equation}
  g \approx 1.6 \times 0.18 \times 0.006 \times P^{0.4}
  \qquad\unit{[m,\,W]}
\end{equation}
For the 30\unit{kW} machine above this gives $g \approx 0.988\unit{mm}$.
Heavy-duty service typically warrants a 60\% increase
($g \approx 1.5\unit{mm}$ here). PM machines often use the smallest
mechanically feasible gap to minimise the required magnet volume.

\subsection{Tip Speed and Centrifugal Stress}

The tangential (tip) speed of the rotor is $v = \omega r$. For a
0.25\unit{m} radius rotor to reach Mach~1
($v \approx 345\unit{m/s}$) the angular speed is
$\omega = v/r \approx 1400\unit{rad/s}$, i.e.\ $n \approx 14\,000\unit{rpm}$ --
not feasible for any reasonable rotor.

The centrifugal radial stress in a smooth solid rotor of density
$\rho$ at the surface is
\begin{equation}
  \sigma_\text{mech} = C'\,\rho\,\omega^2\,r^2,
  \qquad
  C' = \frac{3+\nu}{4}
\end{equation}
where $\nu$ is Poisson's ratio. For steel, $\nu \approx 0.29$, so
$C' \approx 0.823$.

\paragraph{Worked example -- Pyrh\"onen 6.3.}
A smooth steel cylinder with a small bore, $\rho = 7860\unit{kg/m^3}$,
yield strength $\sigma_y = 300\unit{MPa}$, $\nu = 0.29$, spinning at
$n=15\,000\unit{rpm}$.
\begin{equation}
  r_\text{max} = \sqrt{\frac{\sigma_y}{C'\,\rho\,\omega^2}}
              = \sqrt{\frac{300\cdot 10^6}{0.823\cdot 7860\cdot(2\pi\cdot 15000/60)^2}}
              \approx 0.13\unit{m}
\end{equation}
For comparison, a 0.25\unit{m} radius at the same speed would
\emph{exceed} the yield strength -- this is the fundamental reason
that high-speed machines must be slender (and why the aspect ratio
$\chi$ grows for high-speed designs).

\subsection{Slot-Height Optimisation}

Tall slots offer more area for copper, allowing higher electric
loading $\bar A$. But there are tradeoffs:
\begin{itemize}[noitemsep]
  \item Cooling becomes harder (heat has to travel further out of the
        deep slot).
  \item Slot leakage flux grows.
  \item If $D_\text{rotor}$ is held constant, the outer diameter
        grows; if $D_\text{out}$ is held constant, the rotor shrinks
        and the torque drops because $T \propto D^2 L$.
\end{itemize}
A standard optimisation argument shows that the optimum normalised
slot height $d = h_\text{slot}/(R_\text{out}-R_\text{rotor})$ is about
0.5: torque scales as $T \propto (1-d)d$, maximised at $d = 1/2$.
A common practical range is $0.3 \le d \le 0.6$.

% =====================================================================
\section{Main Flux Paths}
\label{sec:carter}

\subsection{Carter's Coefficient}

Stator slotting causes the air-gap flux density to dip above each slot
opening, raising the effective reluctance of the gap. Carter's
coefficient $k_c \ge 1$ replaces the slotted stator with a smooth
stator of enlarged effective gap:
\begin{equation}
  g_\text{eff} = k_c\,g
\end{equation}
The simplified Pyrh\"onen form (Eq.~3.34) is
\begin{equation}
  k_c = \frac{\tau_u}{\tau_u - b_e},
  \quad
  b_e = \kappa\,b_1,
  \quad
  \kappa = \frac{b_1/g}{5 + b_1/g}
  \label{eq:carter}
\end{equation}
with $\tau_u$ the slot pitch at the bore, $b_1$ the slot opening, and
$g$ the physical gap. $\kappa\to 0$ for small $b_1/g$ (no slotting
effect) and $\kappa\to 1$ for very wide open slots.

\paragraph{Worked example -- Pyrh\"onen 3.1.}
$g = 0.8\unit{mm}$, $b_1 = 3\unit{mm}$, $\tau_u = 10\unit{mm}$.
\begin{align*}
  \kappa &= \frac{b_1/g}{5 + b_1/g}
          = \frac{3/0.8}{5 + 3/0.8}
          = \frac{3.75}{8.75} = 0.429\\
  b_e    &= 0.429 \cdot 3 = 1.29\unit{mm}\\
  k_c    &= \frac{10}{10 - 1.29} = 1.148\\
  g_\text{eff} &= 0.8\cdot 1.148 = 0.918\unit{mm}
\end{align*}

\paragraph{Armature reaction.}
In addition to the magnet (or rotor) field, the stator current sheet
creates a field of its own that adds to the no-load gap field. For a
linear core the armature reaction merely shifts the operating point;
for a non-linear core close to saturation the reaction effect can
dominate, distorting the field shape further.

\subsection{Effective Core Length}

The axial extent of a real stator stack is broken up by cooling ducts
between packs of laminations. Flux fringes from one pack to the next
across each duct, so the effective magnetic length $L'$ is somewhat
greater than the iron length and somewhat less than the total stack
length.

\paragraph{Worked example -- Pyrh\"onen 3.2.}
25 stator sectors of 30\unit{mm} each separated by 24 cooling ducts of
$b_v = 10\unit{mm}$ each. The total stack length is
$L = 25\cdot 30 + 24\cdot 10 = 990\unit{mm}$ (the iron length alone is
$750\unit{mm}$). With air gap $g = 3\unit{mm}$:
\begin{equation}
  \kappa_v = \frac{b_v/g}{5 + b_v/g} = \frac{10/3}{5 + 10/3} = 0.4,
  \quad
  b_{ve} = 0.4 \cdot 10 = 4\unit{mm}
\end{equation}
The effective core length is then approximately
\begin{equation}
  L' \approx L - n_v\,b_{ve} + 2g
         = 990 - 24\cdot 4 + 2\cdot 3
         = 900\unit{mm}
\end{equation}
(the $+2g$ term accounts for axial-end fringing).

% =====================================================================
\section{Magnetizing Inductance}
\label{sec:Lm}

The magnetizing inductance $L_m$ of an AC machine is the
self-inductance of one phase due to the fundamental-component
air-gap field it creates when carrying only its own current and no
saturating component. It is the dominant element of the shunt branch
in the per-phase equivalent circuit (Section~\ref{sec:eqckt}).

\subsection{Flux per Pole}

A sinusoidally distributed air-gap flux density of peak value
$\hat B$ integrated over half a wavelength (one pole pitch $\tau_p$,
on axial length $L'$) gives the peak flux per pole:
\begin{equation}
  \hat\Phi_\text{pole}
    = \int_0^\pi \hat B \sin\theta \cdot \frac{\tau_p}{\pi} L'\,d\theta
    = \frac{2}{\pi}\,\hat B\,\tau_p\,L'
  \label{eq:phi_pole}
\end{equation}
The corresponding stator flux linkage is
$\lambda = k_{w,1}\,N_s\,\hat\Phi_\text{pole}$.

\subsection{MMF of a Polyphase Winding}

For a balanced $m$-phase winding with $N_s$ turns per phase, the
fundamental space MMF amplitude is
\begin{equation}
  \widehat{\mathrm{MMF}}_1
    = \frac{4}{\pi}\,\frac{k_{w,1}\,N_s}{2p}\,\sqrt{2}\,I_{s,\text{rms}}
  \label{eq:mmf1}
\end{equation}
The factor $\tfrac{4}{\pi}$ is the Fourier amplitude of a square-wave
ampere-turn distribution; division by $2p$ converts total turns to
turns per pole; the $\sqrt{2}$ converts RMS to peak.

\subsection{$L_m$ Derivation}

Setting the MMF equal to $\hat H\,g_\text{eff} =
\hat B\,g_\text{eff}/\mu_0$ and solving for $\hat B$:
\begin{equation}
  \hat B
    = \frac{\mu_0}{g_\text{eff}}\cdot
      \frac{4}{\pi}\,\frac{k_{w,1}N_s}{2p}\,\sqrt{2}\,I_{s,\text{rms}}
\end{equation}
The flux linkage per phase becomes
\begin{equation}
  \lambda = \frac{2}{\pi}\,\frac{\mu_0}{g_\text{eff}}\,
            \frac{4}{\pi}\,\frac{(k_{w,1}N_s)^2}{2p}\,
            \tau_p\,L'\,\sqrt{2}\,I_{s,\text{rms}}
\end{equation}
Dividing by $\sqrt{2}\,I_{s,\text{rms}}$ gives the per-phase
magnetizing inductance:
\begin{equation}
  L_m^{(\text{ph})}
    = \frac{2\mu_0 D}{p^2\,\pi\,g_\text{eff}}\,L'\,(k_{w,1}N_s)^2
  \label{eq:Lm_phase}
\end{equation}
where $\tau_p = \pi D/(2p)$ has been used. For three-phase analysis
one usually quotes the \emph{cyclic} magnetizing inductance
\begin{equation}
  L_m = \frac{3}{2}\,L_m^{(\text{ph})}
  \label{eq:Lm_cyclic}
\end{equation}
which is the inductance seen in the equivalent dq circuit.

\paragraph{Induction motor implications.}
$L_m \propto 1/g_\text{eff}$. A larger effective gap reduces $L_m$,
increases the magnetising current $I_m$, lowers the power factor, and
ultimately reduces efficiency. This is why induction motors are
designed with the smallest mechanically feasible gap (typically a
fraction of a millimetre for sub-kW motors, a few millimetres at
multi-MW).

% =====================================================================
\section{Leakage Flux}
\label{sec:leakage}

Flux that links \emph{some but not all} of the winding turns is called
leakage flux. It does not contribute to torque, but it does store
energy and shows up in the equivalent circuit as the leakage
reactance.

\subsection{Slot Leakage}

Flux closing around a single tooth slot without linking the rotor
winding constitutes slot leakage. Applying Amp\`ere's law on a closed
contour inside the slot,
\begin{equation}
  \oint \vec H\cdot d\vec\ell = NI,
\end{equation}
the contour can be split into a vertical $H_2$ portion (close to the
slot bottom, with most of the enclosed turns) and a horizontal $H_1$
portion (across the slot mouth, with few turns enclosed).
Since $H_2 \gg H_1$ but $\ell_2 \ll \ell_1$, the dominant contribution
is the slot-bottom path; this is the source of the large slot-leakage
inductance.

\subsection{End-Winding Leakage}

End-windings carry the full phase current but do not contribute to the
useful air-gap field. The associated leakage inductance is roughly
proportional to the end-winding length and is one of the reasons why
fractional-slot tooth-coil windings (with very short end turns) achieve
high torque density.

\subsection{Magnet-to-Magnet and Tooth-Tip Leakage}

In PM machines, a fraction of the magnet flux closes from one magnet
to its neighbour through the air at the magnet ends, and through the
tooth tips back to the stator slots. These leakage paths can be
modelled with a reluctance network and are responsible for the
reduction of the useful air-gap fundamental below the value computed
from a 1D MEC.

\subsection{Cogging Torque}

A static interaction between the magnets and the stator slot openings
produces a position-dependent reluctance, hence a position-dependent
torque even at zero stator current. This \emph{cogging torque} is a
torque ripple, not useful for propulsion. Fractional-slot windings
attenuate it; further reduction comes from skewing the stator or the
magnets by one slot pitch.

% =====================================================================
\section{Thermal Modelling}
\label{sec:thermal}

\subsection{Electrical--Thermal Analogy}

Heat flow through a solid obeys the same diffusion equation as
electric current in a resistor:
\begin{equation}
  q = \frac{\Delta T}{R_\text{th}}
  \qquad\Longleftrightarrow\qquad
  I = \frac{\Delta V}{R}
\end{equation}
The analogy is summarised in Table~\ref{tab:thermalanalogy}.

\begin{table}[H]
  \centering
  \caption{Electrical--thermal analogy.}
  \label{tab:thermalanalogy}
  \begin{tabular}{ll}
    \toprule
    Electrical & Thermal \\
    \midrule
    Voltage $V$\unit{[V]}              & Temperature $T$\unit{[K]} \\
    Current $I$\unit{[A]}              & Heat flow $q$\unit{[W]} \\
    Resistance $R$\unit{[\Omega]}      & Thermal resistance $R_\text{th}$\unit{[K/W]} \\
    Capacitance $C$\unit{[F]}          & Heat capacity $m c_p$\unit{[J/K]} \\
    Conductivity $\sigma$               & Thermal conductivity $k$\unit{[W/(m\!\cdot\!K)]} \\
    \bottomrule
  \end{tabular}
\end{table}

A lumped thermal network of resistors and capacitors -- one node per
component, one resistor per heat-flow path, one capacitor per thermal
mass -- can predict steady-state and transient temperature rise
without solving any PDE.

\subsection{Conduction}

For a slab of thickness $\ell$, area $A$ and conductivity $k$:
\begin{equation}
  R_\text{th,cond} = \frac{\ell}{k\,A}
\end{equation}
A copper slot, for example, conducts heat radially outward through the
slot insulation (low $k$), the stator tooth iron (medium $k$), and
into the frame.

\subsection{Convection}

For a surface of area $A$ exposed to a coolant with film coefficient
$h$:
\begin{equation}
  R_\text{th,conv} = \frac{1}{h\,A}
\end{equation}
Natural convection $h\sim 5$--$25\unit{W/(m^2K)}$;
forced air $h\sim 25$--$250\unit{W/(m^2K)}$;
water cooling $h\sim 10^3$--$10^4\unit{W/(m^2K)}$.

\paragraph{Fan-load scaling.}
For a fan, the torque scales as $T\propto\omega^2$ and the power as
$P=T\omega\propto\omega^3$. Doubling the airflow speed therefore
requires eight times the shaft power -- which is why uprating the
cooling of an existing machine by simply spinning the fan faster has
diminishing returns.

\subsection{Worked Example -- Thermal Uprate}

A motor rated $P_\text{out} = 10\unit{kW}$, $\eta = 0.9$,
$T_\text{amb} = 40\unit{\degree C}$, $T_\text{motor} = 80\unit{\degree C}$
at rating. So $P_\text{loss} = P_\text{out}/\eta - P_\text{out} =
1\unit{kW}$ and $\Delta T = 40\unit{K}$, giving the steady-state
thermal resistance $R_\text{th} = 40\unit{K/kW} = 0.04\unit{K/W}$.

\emph{Uprate by 50\%.} If the same motor must deliver $15\unit{kW}$,
the $I^2R$ losses scale as $I^2 \propto P^2$, so
$P_\text{loss}^{\text{new}} = 1\,\text{kW}\cdot(1.5)^2 = 2.25\unit{kW}$
and the temperature rise becomes $\Delta T = R_\text{th}\cdot
P_\text{loss}^{\text{new}} = 40 \cdot 2.25 = 90\unit{K}$ above
ambient, i.e.\ $T_\text{motor} \approx 130\unit{\degree C}$.
A more careful calculation that also lets the winding resistance grow
with temperature (a $\sim 25$\% increase at $130\unit{\degree C}$
relative to $80\unit{\degree C}$) gives
$\Delta T \approx 112.5\unit{K}$ and $T_\text{motor}\approx 152\unit{\degree C}$ --
likely beyond Class~F insulation limits, so the uprate is not
sustainable without improved cooling.

\subsection{Thermal Network for a Machine}

A practical thermal model groups losses by region (copper, iron,
end winding, magnet) and connects them through the local conduction
and convection resistances. Heat then flows from copper $\to$ slot
insulation $\to$ tooth iron $\to$ back-iron $\to$ frame $\to$ ambient.
Air ducts add parallel convection paths from the iron to the duct air,
in series with the duct-to-ambient resistance. The end-winding region
has its own convection path to the internal air volume.

% =====================================================================
\section{Permanent Magnet Machines}
\label{sec:pm}

\subsection{Hard vs.\ Soft Ferromagnetic Materials}

The B--H loop of a \emph{soft} ferromagnetic material (silicon steel,
nickel-iron alloys) has small enclosed area and small coercive force
$H_c$: it follows the applied field easily and reverses with little
hysteresis loss. Soft materials are used for the stator and rotor
cores of all electrical machines.

A \emph{hard} ferromagnetic material (ferrite, samarium-cobalt,
neodymium-iron-boron) has very large $H_c$ and large remanent flux
density $B_r$. Once magnetised it retains its magnetisation against
substantial demagnetising fields. The figure of merit is the
\emph{maximum energy product} $|BH|_\text{max}$ -- the area of the
largest rectangle that fits in the second-quadrant
$B$--$H$ characteristic.
Typical values:
\begin{itemize}[noitemsep]
  \item Ferrite: $B_r \approx 0.4\unit{T}$, $H_c\approx 250\unit{kA/m}$,
        $|BH|_\text{max}\approx 30\unit{kJ/m^3}$.
  \item NdFeB: $B_r\approx 1.2$--$1.4\unit{T}$,
        $H_c\approx 800$--$1200\unit{kA/m}$,
        $|BH|_\text{max}\approx 250$--$430\unit{kJ/m^3}$.
\end{itemize}

\subsection{Magnet Equivalent Circuit}

In the linear (recoil) region of the second quadrant, the magnet
behaves like an MMF source $H_c\,l_m$ in series with an internal
reluctance $\mathcal{R}_m = l_m/(\mu_0\mu_r A_m)$, where
$\mu_r\approx 1.05$ for NdFeB.
Equivalently, on the $B$ axis, the source is the remanent flux
$\Phi_r = B_r A_m$.
The whole magnetic circuit becomes a flux source driving an external
reluctance (the air gap, the iron, and any leakage paths) in series
with $\mathcal{R}_m$:
\begin{equation}
  \Phi_g = \frac{B_r\,A_m}{\,1 + \mathcal{R}_g/\mathcal{R}_m\,}
\end{equation}
In the idealised 1D case with $A_m = A_g = A$ and infinite iron
permeability, this reduces (cf.\ Section~5 of the HW2 report) to
\begin{equation}
  B_g = \frac{B_r\,l_m}{l_m + \mu_r\,g}
\end{equation}
With Carter's correction, $g$ is replaced by $k_c g$.

\subsection{Operating Point and Sizing Intuition}

The operating point of the magnet is the intersection of its
demagnetisation line with the load line $-B/(\mu_0 H) =
\mathcal{R}_m/\mathcal{R}_g$. Moving the operating point closer to
$|BH|_\text{max}$ -- a load line through the ``knee'' of the
recoil curve -- maximises the energy delivered per unit volume of
magnet material.

\paragraph{Doubling the magnet thickness does not double $B_g$.}
The MMF source doubles but the magnet reluctance also doubles, while
the air-gap reluctance is unchanged. The increase in $B_g$ is
therefore sub-linear: e.g.\ for $l_m = 4\unit{mm}$, $g = 4\unit{mm}$
the gap field might be $B_g = 0.6\unit{T}$; doubling the magnet to
$l_m = 8\unit{mm}$ raises $B_g$ to roughly $0.8\unit{T}$, not
$1.2\unit{T}$.
Conversely, a magnet thinner than $\mu_r g$ is essentially short-circuited
by its own reluctance and delivers a small fraction of $B_r$.

\paragraph{Maximum energy transfer.}
The classical maximum-power-transfer theorem applies: the energy
density delivered to the external circuit is maximised when the
external (load) reluctance equals the internal (magnet) reluctance.
For a fixed magnet, this fixes the optimum effective air gap; for a
fixed air gap, it fixes the optimum magnet thickness:
\begin{equation}
  l_m^\text{opt} \;\approx\; \mu_r\,g_\text{eff}.
\end{equation}
In practice, the working point is pushed slightly toward the
high-$B$ side of this optimum because demagnetisation safety margin
is more valuable than the last few percent of efficiency.

\subsection{Demagnetisation Safety}

The temperature dependence of $B_r$ (typically $-0.1\%/$\degree C for
NdFeB) and of $H_c$ (typically $-0.5\%/$\degree C) shifts the
demagnetisation line toward the origin as the magnet heats up.
A short-circuit fault, an inverter shoot-through, or a starting
current transient can drive the operating point below the
\emph{knee} of the recoil line; below the knee, recoil is irreversible
and the magnet loses some of its $B_r$ permanently. Conservative PM
machine design therefore checks the worst-case operating point
(highest temperature, largest stator current, demagnetising stator
MMF) against a margin to the knee.

% =====================================================================
\appendix
\section{Summary of Key Formulas}

\begin{itemize}[noitemsep]
  \item Slots per pole per phase: $q = N_s/(2pm)$.
  \item Distribution factor:
        $k_{d,n} = \sin(nq\alpha_e/2)/[q\sin(n\alpha_e/2)]$.
  \item Pitch factor: $k_{p,n} = \sin(n\lambda/2)$, with $\lambda$ the
        coil span in electrical degrees.
  \item Winding factor: $k_{w,n} = k_{d,n}\,k_{p,n}$.
  \item Flux per pole: $\hat\Phi_\text{pole} = (2/\pi)\hat B\,\tau_p L'$.
  \item Alternator EMF (per harmonic): $V_\text{rms} = 4.44\,f\,N_s\,k_{w,n}\,\hat B\,A_\text{pole}$.
  \item Output equation: $S_i = (\pi^2/\sqrt{2})\,k_{w,1}\,\bar A\,\hat B\,D^2 L\,n_\text{syn}
                              = C\,D^2 L\,n_\text{syn}$.
  \item Carter's coefficient: $k_c = \tau_u/(\tau_u-b_e)$,
        $b_e = \kappa b_1$, $\kappa = (b_1/g)/(5+b_1/g)$.
  \item Effective core length: $L' \approx L - n_v b_{ve} + 2g$.
  \item Magnetising inductance per phase:
        $L_m^{(\text{ph})} = (2\mu_0 D/(p^2\pi g_\text{eff}))\,L'\,(k_{w,1}N_s)^2$.
  \item Shear stress / Maxwell stress: $\sigma = B^2/(2\mu_0)$.
  \item Air-gap flux density of a surface-PM machine:
        $B_g = B_r l_m/(l_m + \mu_r k_c g)$.
\end{itemize}

\section{Symbols and Units}

\begin{table}[H]
  \centering
  \begin{tabular}{lll}
    \toprule
    Symbol & Meaning & Typical unit \\
    \midrule
    $D$         & stator bore diameter         & m \\
    $L$         & axial (stack) length         & m \\
    $L'$        & effective magnetic length    & m \\
    $g$         & physical air-gap length      & mm \\
    $g_\text{eff}$ & effective air-gap (with $k_c$) & mm \\
    $\tau_s$    & slot pitch at the bore       & mm \\
    $\tau_p$    & pole pitch at the bore       & mm \\
    $b_1, w_s$  & slot opening width           & mm \\
    $w_{tb}$    & tooth body width             & mm \\
    $h_\text{slot}$ & slot height               & mm \\
    $N_s$       & total slot count             & --- \\
    $2p$        & pole count                   & --- \\
    $m$         & phase count                  & --- \\
    $q$         & slots per pole per phase     & --- \\
    $\alpha_e$  & slot pitch, electrical       & deg \\
    $k_d$, $k_p$, $k_w$ & winding factors       & --- \\
    $\bar B$    & specific magnetic loading    & T \\
    $\bar A$    & specific electric loading    & A/m \\
    $C$         & specific machine constant    & W$\cdot$s/m\textsuperscript{3} \\
    $B_r$       & magnet remanence             & T \\
    $H_c$       & magnet coercivity            & A/m \\
    $\mu_r$     & magnet relative permeability & --- (usually $\sim 1.05$) \\
    $l_m$       & magnet radial thickness      & mm \\
    \bottomrule
  \end{tabular}
\end{table}

\end{document}
